% =============================================================================
% Semana 09 — Design de Informação: Pré-atenção, Hierarquia Visual e Chartjunk
% Aula de 3 horas (19h–22h) | 29/04/2026
% =============================================================================

\chapter{Semana 09 (29/04/2026) — Design de Informação: do Visual Correto ao Visual Claro}

% ---------------------------------------------------------------------------
\section{Objetivo da semana}
% ---------------------------------------------------------------------------

Pessoal, um gráfico pode estar tecnicamente correto e ainda assim falhar
completamente na comunicação. O leitor não consegue encontrar o que importa,
os eixos são confusos, as cores não têm propósito. Hoje aprendemos a
transformar visuais corretos em visuais \textbf{claros} — onde a mensagem é
imediata, a hierarquia é evidente e cada elemento visual ganha seu espaço
por um motivo.

\begin{BoardBox}
\textbf{Roteiro da aula (19h–22h)}
\begin{enumerate}
  \item 19h–19h30 — \textbf{TEORIA:} Pré-atenção — como o cérebro lê um gráfico
  \item 19h30–20h — \textbf{TEORIA:} Teste do Relance + Título que é mensagem
  \item 20h–20h20 — \textbf{TEORIA:} Chartjunk (Tufte) + Hierarquia visual
  \item 20h20–20h40 — \textbf{DATASET:} Importar Video Game Sales no Tableau
  \item 20h40–21h30 — \textbf{PRÁTICA GUIADA:} Dashboard de vendas de games — linha, scatter, heatmap
  \item 21h30–21h50 — \textbf{PRÁTICA:} Antes e Depois — aplicar design no próprio dashboard
  \item 21h50–22h — \textbf{Revisão do T5 + revisão Q2}
\end{enumerate}
\end{BoardBox}

\begin{NoteBox}
\textbf{Leitura base:} \textit{Datastory}, Nancy Duarte — Cap.\ 4 e 5 (design de slides e hierarquia visual).\\
\textbf{Leitura complementar:} \textit{The Visual Display of Quantitative Information}, Edward Tufte — conceito de data-ink ratio e chartjunk.
\end{NoteBox}

% =============================================================================
\section{Parte I — Pré-atenção (19h–19h30)}
% =============================================================================

\begin{SolvedBox}
\textbf{O que é pré-atenção?}

Antes de ler deliberadamente qualquer coisa em um gráfico, o cérebro
humano \emph{automaticamente} detecta certos atributos visuais em menos
de 250 milissegundos — sem esforço consciente. Esses são os atributos
\textbf{pré-atencionais}.

Se você usa esses atributos para destacar o que importa, o leitor chega
à mensagem antes mesmo de ler o título. Se você os usa aleatoriamente
(tudo em cores diferentes, tamanhos mistos), o leitor fica desorientado.
\end{SolvedBox}

\begin{BoardBox}
\textbf{Atributos pré-atencionais — os mais úteis em visualização de dados}

\begin{tabular}{|l|p{5.5cm}|p{5cm}|}
\hline
\textbf{Atributo} & \textbf{O que faz} & \textbf{Uso recomendado} \\
\hline
\textbf{Cor} & Agrupa, destaca e emociona & Usar 1 cor de destaque; resto em cinza/neutro \\
\textbf{Tamanho} & Indica magnitude & Barra maior = valor maior; consistência obrigatória \\
\textbf{Posição} & Organiza hierarquia & Elemento mais importante no canto superior esquerdo \\
\textbf{Intensidade} & Contraste claro-escuro & Tom mais escuro = mais importante \\
\textbf{Forma} & Diferencia categorias & Formas distintas (círculo vs quadrado) para subgrupos \\
\textbf{Enclosure} & Agrupa com borda ou área & Caixa ao redor de um grupo de dados relacionados \\
\hline
\end{tabular}

\textbf{Regra de ouro:} use, no máximo, \textbf{2 atributos pré-atencionais} por visual.
Mais do que isso e a hierarquia colapsa — o olho não sabe para onde ir.
\end{BoardBox}

\begin{SolvedBox}
\textbf{Princípio da cor com propósito:}

\begin{itemize}
  \item \textbf{Cor semântica:} vermelho = alerta/negativo; verde = positivo/ok;
  amarelo = atenção. Use onde há convenção estabelecida — não inverta sem motivo.
  \item \textbf{Cor de destaque:} 1 cor viva para o elemento que carrega a mensagem;
  tudo mais em cinza. O contraste leva o olho diretamente ao dado importante.
  \item \textbf{Cor para categorias:} máximo de 6 cores distintas. Acima disso, o leitor
  perde a associação cor-categoria.
  \item \textbf{Acessibilidade:} 8\% dos homens têm daltonismo vermelho-verde.
  Evite depender exclusivamente dessa combinação como único diferenciador.
\end{itemize}
\end{SolvedBox}

% =============================================================================
\section{Parte II — Teste do Relance e Título-Mensagem (19h30–20h)}
% =============================================================================

\begin{BoardBox}
\textbf{O Teste do Relance (3 segundos)}

Mostre seu gráfico por 3 segundos e cubra. Pergunte à audiência:
\begin{enumerate}
  \item O que esse gráfico está mostrando?
  \item Qual é a mensagem principal?
  \item O que você deveria fazer com essa informação?
\end{enumerate}

Se a audiência não conseguir responder à Q2 em 3 segundos, o visual falhou —
não o leitor. O problema está no design.

\textbf{Diagnóstico do fracasso no Teste do Relance:}
\begin{itemize}
  \item Título descritivo/neutro: ``Dados de Vendas 2025'' — não passa informação.
  \item Excesso de elementos: o olho não sabe onde pousar primeiro.
  \item Cores sem propósito: todas as barras em cores diferentes sem razão.
  \item Legenda separada: o leitor precisa desviar o olhar para decodificar.
\end{itemize}
\end{BoardBox}

\begin{SolvedBox}
\textbf{Títulos rótulo vs títulos mensagem:}

\begin{tabular}{|l|l|}
\hline
\textbf{Título rótulo (ruim)} & \textbf{Título mensagem (bom)} \\
\hline
``Vendas por Região'' & ``Região Sul concentra 58\% das vendas'' \\
``Evolução Mensal'' & ``Churn caiu 34\% após o programa de fidelidade'' \\
``Satisfação de Clientes'' & ``NPS Premium é 2× superior ao do plano Básico'' \\
``Comparativo de Custos'' & ``Canal físico custa R\$12 por venda; digital, R\$0,80'' \\
\hline
\end{tabular}

\textbf{Fórmula do título mensagem:} [Sujeito do dado] + [verbo de tendência/comparação]
+ [magnitude/contexto]. Ex.: ``Retenção \textbf{aumentou} 22pp com o app'' —
sujeito (retenção) + verbo (aumentou) + magnitude (22pp).
\end{SolvedBox}

% =============================================================================
\section{Parte III — Chartjunk + Hierarquia Visual (20h–20h20)}
% =============================================================================

\begin{BoardBox}
\textbf{Data-ink ratio (Tufte) — para o quadro}

Edward Tufte (1983) propõe o conceito de \textbf{data-ink ratio}:

\[
\text{Data-ink ratio} = \frac{\text{Tinta usada para representar dados}}{\text{Tinta total do gráfico}}
\]

Maximize esse ratio. Todo elemento visual que não carrega dado é candidato
à remoção.
\end{BoardBox}

\begin{SolvedBox}
\textbf{Exemplos de Chartjunk — o que remover:}

\begin{tabular}{|l|p{8cm}|}
\hline
\textbf{Elemento} & \textbf{Por que remover?} \\
\hline
Grades pesadas & O olho vai para as linhas, não para os dados \\
3D em barras/pizzas & Distorce o tamanho percebido; pizza 3D é enganosa \\
Sombras e reflexos & Decoração — não carrega informação \\
Eixo Y com zero forçado & Às vezes necessário; às vezes esconde variação real \\
Legendas redundantes & Se cada barra já tem rótulo, a legenda é ruim \\
Ícones temáticos & Uma pizza com ícone de comida não passa mais informação \\
Gradientes de cor & Cor sem propósito semântico confunde \\
Mais de 3 casas decimais & Precisão falsa — dados de negócio raramente justificam \\
\hline
\end{tabular}

\textbf{Princípio:} duvide de cada elemento. Pergunte: ``Se eu remover isso,
o gráfico fica menos claro?'' Se a resposta for não, remova.
\end{SolvedBox}

% =============================================================================
\section{Hierarquia Visual (complemento — para o quadro)}
% =============================================================================

\begin{BoardBox}
\textbf{Os quatro pilares da hierarquia visual}

\textbf{1. Tamanho:} o elemento mais importante deve ser o maior.
Título $>$ subtítulo $>$ rótulo de dado $>$ nota de rodapé.

\textbf{2. Peso (negrito):} use negrito para o número que carrega a mensagem.

\textbf{3. Cor:} 1 cor de destaque; resto em cinza ou neutro.

\textbf{4. Posição:} leitura ocidental em Z (canto superior esquerdo primeiro).
\end{BoardBox}

\begin{SolvedBox}
\textbf{Consistência — o princípio silencioso:}

\begin{itemize}
  \item \textbf{Eixos:} padronize o eixo Y para comparáveis.
  \item \textbf{Cores:} Região Sul é sempre azul em todas as telas.
  \item \textbf{Fontes:} máximo de 2 famílias tipográficas.
  \item \textbf{Animações:} só se adicionarem informação.
\end{itemize}
\end{SolvedBox}

% =============================================================================
\section{Dataset da Semana — Video Game Sales (Kaggle)}
% =============================================================================

\begin{BoardBox}
\textbf{Base de dados da semana:} \textit{Video Game Sales}

\textbf{Link para download:}
\url{https://www.kaggle.com/datasets/gregorut/videogamesales}

\textbf{Sobre o dataset:}
\begin{itemize}
  \item $\approx$ 16.500 jogos com vendas globais por região
  \item \textbf{Colunas principais:} Name, Platform, Year\_of\_Release,
  Genre, Publisher, NA\_Sales, EU\_Sales, JP\_Sales, Other\_Sales,
  Global\_Sales
  \item Dados limpos, sem necessidade de joins — pronto para importar
  \item Rica variedade: dimensões categóricas (Genre, Platform) +
  medidas contínuas (vendas) + série temporal (Year)
\end{itemize}

\textbf{Pergunta de negócio da semana:}
\textit{``Qual gênero de jogo e qual plataforma devemos priorizar
para o próximo lançamento?''}
\end{BoardBox}

\begin{NoteBox}
\textbf{Importação no Tableau:} File $\to$ Open $\to$ Text file $\to$
\texttt{vgsales.csv}. Verificar se \texttt{Year} é reconhecido como
número (converter para inteiro se necessário).
Filtrar \texttt{Year} $\neq$ N/A para remover registros sem ano.
\end{NoteBox}

% =============================================================================
\section{Parte IV — Construção Guiada no Tableau (20h20–21h30)}
% =============================================================================

\begin{SolvedBox}
\textbf{Construção guiada — Dashboard ``Panorama da Indústria de Games''}

Vamos construir juntos, aplicando os princípios de design aprendidos.
Cada visual será criado primeiro ``feio'' e depois refinado.

\textbf{Sheet 1 — Linha Temporal: Evolução de Vendas Globais por Ano (15 min):}
\begin{enumerate}
  \item Arrastar \texttt{Year} para Columns (contínuo),
  \texttt{Global\_Sales} (SUM) para Rows
  \item \textbf{Versão ruim:} deixar padrão do Tableau —
  título ``Sheet 1'', cor padrão, sem anotação
  \item \textbf{Aplicar design:}
  \begin{itemize}
    \item Título-mensagem: ``Pico de vendas em 2008 — mercado em queda desde então''
    \item Anotar o ponto de pico (2008) com marcador
    \item Remover grade pesada (Format $\to$ Lines $\to$ Grid Lines = None)
    \item 1 cor sólida (azul escuro) para a linha
  \end{itemize}
  \item \textbf{Teste do Relance:} cubra e pergunte ao colega
\end{enumerate}

\textbf{Sheet 2 — Scatter: NA Sales vs EU Sales por Gênero (15 min):}
\begin{enumerate}
  \item Arrastar \texttt{NA\_Sales} (SUM) para Columns,
  \texttt{EU\_Sales} (SUM) para Rows, \texttt{Genre} para Detail + Color
  \item \textbf{Versão ruim:} 12 gêneros × 12 cores — caos visual
  \item \textbf{Aplicar design:}
  \begin{itemize}
    \item Reduzir: cor cinza para todos, \textbf{destaque em azul}
    apenas para o gênero líder (Action)
    \item Adicionar trend line (Analytics $\to$ Trend Line)
    \item Título: ``Action domina: vendas acima de qualquer outro gênero
    tanto na América do Norte quanto na Europa''
    \item Rótulo direto no ponto de destaque (elimina legenda)
  \end{itemize}
\end{enumerate}

\textbf{Sheet 3 — Heatmap: Gênero × Plataforma (15 min):}
\begin{enumerate}
  \item Arrastar \texttt{Genre} para Rows, \texttt{Platform} para
  Columns, \texttt{Global\_Sales} (SUM) para Color
  \item Tipo: Square (Show Me $\to$ Heat Map)
  \item Paleta: sequencial (branco $\to$ azul escuro)
  \item Filtrar apenas Top 10 plataformas por vendas (senão fica denso)
  \item \textbf{Título:} ``Action + PS2 é a combinação mais vendida da história''
  \item Legenda da escala com unidade: ``Vendas em milhões de USD''
\end{enumerate}

\textbf{Montar o Dashboard (5 min):}
\begin{enumerate}
  \item Linha temporal no topo (contexto temporal)
  \item Scatter + Heatmap lado a lado embaixo
  \item Título geral como Big Idea do dashboard
  \item Paleta consistente: mesma família de azul em todas as sheets
\end{enumerate}
\end{SolvedBox}

\begin{BoardBox}
\textbf{Estudo Estatístico com Video Game Sales (para o quadro)}

\begin{enumerate}
  \item \textbf{Tendência temporal:} olhando a linha temporal, em que
  ano as vendas atingiram o pico? Qual a taxa de queda desde então?
  Adicione uma trend line e observe o $R^2$.

  \item \textbf{Correlação regional:} no scatter NA vs EU, a
  correlação é forte? Se um gênero vende bem na América do Norte,
  vende bem na Europa também? E no Japão? (Troque EU por JP e compare.)

  \item \textbf{Top-N:} crie um ranking das 5 plataformas com mais
  vendas globais. A plataforma líder tem quantas vezes mais vendas
  que a 5ª? Essa concentração é surpreendente?

  \item \textbf{Comparação de distribuição:} compare a média e a
  mediana de \texttt{Global\_Sales} por Genre. Gêneros com média
  muito acima da mediana indicam presença de poucos jogos com
  vendas gigantes (outliers positivos). Quais gêneros têm esse padrão?
\end{enumerate}
\end{BoardBox}

\begin{SolvedBox}
\textbf{Questão rápida — Redesign de Dashboard}

Um estagiário apresentou um dashboard com:
\begin{itemize}
  \item Título: ``Dashboard de Vendas''
  \item 10 gráficos, 6 cores sem padrão, sem espaço em branco
\end{itemize}

Liste \textbf{4 ações de redesign}.

\textbf{Resolução:}
\begin{enumerate}
  \item Título-mensagem: ``Vendas Q1: meta atingida em Eletrônicos,
  15\% abaixo em Vestuário''.
  \item Reduzir para 3–4 visuais prioritários.
  \item Paleta restrita: cinza + verde/vermelho para status.
  \item Hierarquia: KPIs em fonte grande, espaço em branco para separar.
\end{enumerate}
\end{SolvedBox}

% =============================================================================
\section{Parte V — Antes e Depois: Refinar o Dashboard (21h30–21h50)}
% =============================================================================

\begin{BoardBox}
\textbf{Exercício Antes e Depois — no seu próprio dashboard (em dupla)}

Voltem ao dashboard de Video Game Sales que acabaram de construir. Para cada
sheet, apliquem o checklist:

\begin{tabular}{|p{5cm}|l|l|}
\hline
\textbf{Critério} & \textbf{Antes} & \textbf{Depois} \\
\hline
Título é mensagem (não rótulo)? & & \\
Passa no Teste do Relance (3s)? & & \\
No máximo 2 atributos pré-atencionais? & & \\
Chartjunk removido (grade, 3D, sombras)? & & \\
1 cor de destaque + resto em cinza? & & \\
Rótulos diretos (sem legenda separada)? & & \\
\hline
\end{tabular}

\textbf{Tarefa:} ajustem pelo menos 2 sheets aplicando os princípios.
Capturem print do ``antes'' e do ``depois'' para entregar.
\end{BoardBox}

% =============================================================================
\section{Complemento — Integração Python $\to$ Tableau}
% =============================================================================

\begin{BoardBox}
\textbf{Pipeline Python $\to$ Tableau (explicação rápida — para o quadro)}

Na Unidade I vocês usaram Python para AED (pandas, matplotlib).
Na Unidade II vocês usam Tableau para visualização e storytelling.
Como conectar os dois?

\textbf{Estratégia recomendada (sem TabPy):}
\begin{enumerate}
  \item \textbf{Python faz o trabalho pesado:} limpeza, feature engineering,
  cálculos estatísticos, agregações, modelos preditivos.
  \item \textbf{Exportar CSV ou Parquet:} o resultado final do Python
  vira um arquivo limpo que o Tableau consome.
  \item \textbf{Tableau faz a visualização:} importa o arquivo,
  cria dashboards interativos, aplica design e storytelling.
\end{enumerate}

\textbf{Exemplo de pipeline:}
\begin{itemize}
  \item Python: \texttt{df.groupby('Genre')['Global\_Sales'].agg(['sum','mean','median']).to\_csv('genre\_stats.csv')}
  \item Tableau: importar \texttt{genre\_stats.csv} $\to$ barras ordenadas + sparklines
\end{itemize}

\textbf{Estratégia avançada (TabPy):}
\begin{itemize}
  \item \texttt{SCRIPT\_REAL} no Tableau chama código Python em tempo real
  \item Útil para scoring de modelos já treinados (ex.: previsão de churn)
  \item Requer TabPy server rodando — não necessário nesta disciplina,
  mas saibam que existe
\end{itemize}

\textbf{Leitura complementar:} Capítulo 7 do livro (Insights com Tableau e Python).
\end{BoardBox}

% =============================================================================
\section{Parte VI — Revisão + Prova II Q2 (21h50–22h)}
% =============================================================================

\begin{NoteBox}
Hoje revisamos o conteúdo da \textbf{Questão 2} da Prova II:
design de dashboards, Teste do Relance e hierarquia visual.
\end{NoteBox}

\begin{SolvedBox}
\textbf{Cenário (estilo Q2 da Prova II):}

Pedro está criando um dashboard executivo no Tableau para o diretor
financeiro de uma distribuidora. Ele tem duas opções:

\textbf{Opção A:} Título ``Vendas Q4: Meta superada em 12\%'';
KPI grande em verde; 2--3 cores; espaço em branco.

\textbf{Opção B:} Título ``Dashboard de Vendas''; 12 gráficos;
8 cores; sem espaço livre.

\textbf{Questão:} qual é superior e por quê?

\textbf{Resolução:} A é superior — título é mensagem, uso restrito
de cores direciona atenção, espaço em branco organiza, e o dashboard
passa no Teste do Relance.
\end{SolvedBox}

% =============================================================================
\section{Revisão Semanal — Questão tipo Prova II (Q2)}
% =============================================================================

\begin{NoteBox}
Hoje revisamos o conteúdo da \textbf{Questão 2} da Prova II:
design de informação, hierarquia visual e o Teste do Relance.
\end{NoteBox}

\begin{SolvedBox}
\textbf{Cenário (estilo Q2 da Prova II):}

Um estagiário criou um dashboard para o diretor de logística de uma
transportadora. O dashboard contém:
\begin{itemize}
  \item 14 gráficos (pizza, barras, linhas, área, waffle, donut, treemap\ldots)
  \item Todas as 8 cores da paleta padrão do Tableau + gradientes
  \item Títulos genéricos: ``Gráfico 1'', ``Gráfico 2'', etc.
  \item Sem espaço em branco — ``aproveitando toda a tela''
  \item Sem destaque visual para KPIs críticos
\end{itemize}

Você foi chamado para refazer o dashboard aplicando os princípios de
hierarquia visual. Descreva \textbf{4 ações concretas} que faria.

\textbf{Resolução:}
\begin{enumerate}
  \item \textbf{Reduzir para 3--4 visualizações} — selecionar apenas os
  gráficos que respondem à pergunta de decisão do diretor (ex.: custo por
  km, prazo médio de entrega vs.\ meta, entregas atrasadas por rota).
  \item \textbf{Títulos-mensagem} — trocar ``Gráfico 1'' por títulos que
  comunicam o insight: ``Rota SP--RJ concentra 40\% dos atrasos''.
  \item \textbf{Restringir paleta a 2--3 cores} — cinza neutro para a base,
  vermelho para alerta (abaixo da meta), verde para aprovação.
  \item \textbf{KPI cards em destaque} — colocar os indicadores-chave
  (custo/km, prazo médio, \% atrasos) no topo do dashboard com tamanho
  grande, aplicando o Teste do Relance: em 3 segundos o diretor vê se
  está no verde ou no vermelho.
\end{enumerate}

\textbf{Lição:} menos é mais. O objetivo do dashboard executivo é
\emph{acelerar a decisão}, não impressionar com quantidade de gráficos.
\end{SolvedBox}

% ---------------------------------------------------------------------------
\section{Materiais Preparados}
% ---------------------------------------------------------------------------
\begin{itemize}
  \item Slides com galeria de ``Antes e Depois'' (8 exemplos reais)
  \item Handout: checklist de design de informação (10 itens)
  \item Folha de atividade individual impressa com o cenário
  \item Link: Viz of the Day — Tableau Public (exemplos de excelente design)
\end{itemize}

% ---------------------------------------------------------------------------
\section{Próximo passo}
% ---------------------------------------------------------------------------

Na \textbf{Semana 10} chegam as apresentações do Trabalho 5 e aprofundamos
o Tableau: filtros contextuais, parâmetros e ações de dashboard. Vamos
construir dashboards que não apenas mostram dados — mas que forçam a
interação e levam à decisão.

Perguntem aos alunos: \textit{``Pensem no último relatório ou dashboard que
vocês receberam no trabalho ou estágio. Ele passaria no Teste do Relance?
O título dizia a mensagem ou era apenas um rótulo? O que vocês fariam
diferente agora?''}
